% Add the approved author list, affiliations, and correspondence here.
% Author identities remain unfilled. The final abstract is grounded in the locked results.
\newcommand{\manuscripttitle}{TrajRisk: Treatment-Aware Multi-Scale Graph Neural Networks for Trajectory-Informed Risk Propagation in Single-Cell Transcriptomics}
\newcommand{\manuscriptauthors}{}
\newcommand{\manuscriptaffiliations}{}
\newcommand{\manuscriptcorrespondence}{}
\newcommand{\manuscriptabstract}{Single-cell profiles of acetaminophen liver injury reveal heterogeneous hepatocyte responses, but connecting these states across injury and recovery remains challenging. We present TrajRisk, which uses sampling time and treatment annotations to learn directed transition graphs across four temporal scales and integrates expression-derived features through a graph neural network. Its dimensionless score estimates normalized Monocle 3 pseudotime in a mouse injury--recovery reference, rather than adverse-outcome probability. We evaluated TrajRisk across 36,067 mouse hepatocytes from four studies. Within the reference, TrajRisk identified elevated-pseudotime cells at 24~h with higher average precision than expression-only XGBoost (0.498 versus 0.379), although XGBoost performed better overall and at later recovery timepoints. External mouse performance varied by cohort. After mouse-only supervised training in shared mouse--human coordinates, scores distinguished donors with acetaminophen-associated acute liver failure from healthy donors (AUROC 0.844; 19 donors). Spatial signature enrichment occurred at mouse injury boundaries but not in independent pericentral zonation analyses. Graph flux analysis identified a candidate 24--48~h recovery gateway linking compensation-associated states to recovery-proximal states. TrajRisk organizes injury states across time and prioritizes transition-associated programs for experimental testing.}
